\chapter{Construirea setului de date}
\label{section:building}
Setul de date final conține 14.500 de fișiere audio distincte, împărțite pe 46 de genuri muzicale.
Dacă inițial numărul de fișiere era egal pentru fiecare gen, în urma curățării setului de date,
distribuția nu a mai rămas egală, acest aspect fiind tratat în subcapitolul \ref{section:genuri}
Distribuția menționată poate fi observată în figura \ref{fig:f1}.

\begin{figure}[h]
\includegraphics[width=0.9\textwidth]{fisiere.png}
\centering
\captionsetup{justification=centering,font=small,labelfont=small}
\caption{Distribuția fișierelor per genuri muzicale}
\label{fig:f1}
\end{figure}

În continuare sunt prezentate, în mai multe detalii, etapele construirii setului de date.

\section{Colectarea fișierelor audio}
Pentru colectarea fișierelor audio, a fost utilizat API-ul pus la dispoziție de Spotify.
Printr-un request în care este specificat genul dorit, se obține o listă de albume care se încadrează în acest gen.
Pentru fiecare album, se face un request prin care, specificând id-ul albumului se obține lista de cântece din acel album.
Astfel, avem acces la id-ul fiecărei melodii, genul în care se încadrează și un url de vizualizare de unde poate fi descărcat
un segment de 30 de secunde din melodie. Pentru a reține detaliile despre melodiile deja colectate, gen, id-ul melodiei, id-ul
albumului din care provine și url-ul pentru previzualizare, s-a utilizat un cluster de MongoDB. Aceste detalii adiționale vor fi utilizate
ulterior pentru realizarea experimentelor.

Odată ce aceste informații au fost obținute, pentru fiecare melodie care nu este deja înregistrată în baza de date, segmentul
audio va fi descărcat și procesat.

\section{Izolarea elementelor vocale}
Pentru a izola elementele vocale de partea instrumentală am utilizat Spleeter,
un proiect de separare a surselor audio lansat de Deezer în 2019 \footnote{https://deezer.io/releasing-spleeter-deezer-r-d-source-separation-engine-2b88985e797e},
în varianta de API python, pentru a integra ușor acest pas
direct in pipeline-ul de construire a setului de date. Astfel, imediat după descărcare, melodia este automat împărțită în două fișiere, voce
și instrumental.

Pentru a realiza acest lucru, Spleeter utilizează o rețea U-Net \cite{unet}, un auto-encoder convolutional cu conexiuni de tip skip, modificată
pentru a procesa date de tip 3D (timp, frecvență și canale audio), preantrenată pe un set de date de aproape două sute de ore de
înregistrări.\cite{spleeter-paper}

\section{Curățarea setului de date}
În urma separării audio, au fost necesari pași adiționali de procesare a datelor deoarece au fost identificate două probleme.

Prima problemă constă în pauzele vocale din melodiile descărcate ce în urma separării au devenit segmente de liniște care nu dețin
informații în vederea clasificării, reprezintă elemente de zgomot în antrenarea modelului și, așadar, este necesară eliminarea lor.

A doua problemă constă în elemente de instrumental care nu au fost bine separate de modelul preantrenat utilizat.
Problema apare preponderent în cazul genurilor \textit{hard-rock}, \textit{heavy metal}, \textit{death-metal} unde apar solo uri de chitară electrică, iar notele
înalte ale chitarei nu sunt distinse de notele înalte ale artistului,
dar și în cazul melodiilor care conțin instrumente de suflat precum
saxofonul deoarece are spectrul sonor mai apropiat de cel al vocii umane, după cum se poate observa în figura \ref{fig:f2}.

În continuare sunt descriși pașii implementați pentru a rezolva aceste probleme.

\begin{figure}[h]
\includegraphics[width=0.85\textwidth]{comparatie e4.png}
\centering
\captionsetup{justification=centering,font=small,labelfont=small}
\caption[Caption for LOF]{Comparație a spectrogramelor pentru nota E4 obținute cu \textit{ChromeMusicLab} \protect\footnotemark[1]}
\label{fig:f2}
\end{figure}

\footnotetext[1]{\url{https://musiclab.chromeexperiments.com/Spectrogram/}}

\subsection{Eliminarea porțiunilor de liniște}
Pentru a rezolva prima problemă am folosit funcția \textit{librosa.effects.split}
\footnote{\href{http://man.hubwiz.com/docset/LibROSA.docset/Contents/Resources/Documents/generated/librosa.effects.split.html}{librosa.effects.split}}.
Această funcție separă semnalul audio în segmente care nu conțin liniște. După apelarea funcției, segmentele rezultate vor fi concatenate,
rezultând un segment audio fără pauze de liniște după cum poate fi observat în figura \ref{fig:f3}

\begin{figure}[h]
\includegraphics[width=0.5\textwidth]{cleaned_vocals.png}
\centering
\captionsetup{justification=centering,font=small,labelfont=small}
\caption{Waveform înainte și după curățarea porțiunilor de liniște}
\label{fig:f3}
\end{figure}

Noțiunea de liniște în acest caz este considerată ca fiind o secvență cu un număr de decibeli mai mic decât un threshold
(threshold-ul utilizat a fost 35 de decibeli) sub numărul de decibeli de referință. Numărul de decibeli de referință este numărul
maxim de decibeli din semnalul sonor.
Această procesare a fost integrată în pipeline-ul de construire a setului de date, curățarea segmentelor de liniște fiind realizată imediat
după separarea segmentului inițial în voce și instrumental. De asemenea, dacă după eliminarea liniștei semnalul audio rezultat nu conține
cel puțin numărul de secunde per segment necesar la procesare, melodia nu va fi inclusă în setul de date deoarece nu oferă destule
informații vocale.

\subsection{Identificarea și eliminarea fișierelor care conțin doar liniște}
A fost menționat faptul că numărul de decibeli de referință este numărul maxim de decibeli din semnalul sonor. Acest aspect este important
deoarece, în cazul segmentelor care după separare conțin doar liniște, numărul de decibeli de referință este mic iar nicio valoare nu va
fi mai mică decât un threshold sub numărul de referință.

O consecință a acestui aspect este faptul că nu se va identifica nicio porțiune de
liniște. Pentru a rezolva acest caz, se calculează media pătratică a waveform-ului (RMS) pentru fiecare frame și se compară această valoare cu un
threshold (0.05). O reprezentare a valorilor mediilor pătratice pentru frame uri în comparație cu valorile waveform ului poate fi observată
în figura \ref{fig:f4}

\begin{figure}[h]
\includegraphics[width=0.7\textwidth]{rms normal.png}
\centering
\captionsetup{justification=centering,font=small,labelfont=small}
\caption{Valorile RMS}
\label{fig:f4}
\end{figure}

RMS a fost ales în favoarea amplitudinii maxime deoarece evidențiază intensitatea semnalului audio pentru o perioadă mai lungă de timp.
Valoarea amplitudinii maxime pentru un punct în timp izolat nu este reprezentivă pentru întregul frame și de aceea s-a optat pentru
media amplitudinilor maxime, modul în care acestea evoluează fiind și mai apropiat de modul în care oamenii percep schimbările de intensitate.
Diferența poate fi observată în figura \ref{fig:f5}

\begin{figure}[h]
\includegraphics[width=0.9\textwidth]{rms_vs_amp.png}
\centering
\captionsetup{justification=centering,font=small,labelfont=small}
\caption{Valorile RMS în comparație cu amplitudinea maximă}
\label{fig:f5}
\end{figure}

\subsection{Identificarea și eliminarea segmentelor care conțin doar zgomote}
În cazul elementelor de instrumental care nu au fost bine separate de modelul preantrenat utilizat, chiar dacă nu sunt prezente
elemente de voce, segmentul nu este identificat ca unul ce conține doar liniște deoarece instrumentalul care nu a fost separat cu
succes are o intensitate destul de mare. Un segment care se încadrează în această categorie poate fi ascultat \href{https://soundcloud.com/roxana-irimia-480512823/separated-track-that-contains-noise-from-death-angels-mistress-of-pain/s-KWqrjvzkfeT}{aici}.

Pentru a identifica acest tip de segmente s-a ținut cont de faptul că pentru zgomote, perioada de traversare la nivel de zero
(\textit{Zero Crossing
Rate}) este mai mare decât cea a vocii umane. Așadar se poate deduce că, în cazul unui segment care conține doar zgomote, dacă valoarea
maximă pentru perioada de traversare la nivel de zero este mai mare decât un anumit threshold (0.5), iar valoarea RMS maximă este mai mică decât un
anumit threshold (0.4), atunci conține doar zgomote și va fi eliminat din setul de date.
Comparația dintre un segment care conține elemente vocale și un segment care conține doar zgomote poate fi observată în figura \ref{fig:f6}

\begin{figure}[h]%
\centering
\captionsetup{justification=centering,font=small,labelfont=small}
\subfloat[\centering RMS și ZCR pentru voce]{{\includegraphics[width=0.465\textwidth]{rms_zcr_normal.png} }}%
\qquad
\subfloat[\centering RMS și ZCR pentru zgomot]{{\includegraphics[width=0.465\textwidth]{rms_zcr.png} }}%
\caption{Comparație între RMS și ZCR pentru voce și zgomot}%
\label{fig:f6}%
\end{figure}
